\section{The scalar residue on a two-letter arc ending in a down letter}

% BEGIN D83L1
\paragraph{D83L1. Removing the unknown power from the down-letter residue.}
In U83M1's two-letter arc ending in a down letter let
$y=s_{c-1}\ge1$ be the size at its initial up letter and let
$B=s_c+S$. At this component write $P_0=v+p$, $P_1=\varphi P_0$.
The native scalar hypotheses of $\mathrm{InFF}$ say exactly that
$f=F/P_1^y$ and, for some unknown $k\ge0$ and $x_0\in O$,
\[
                 p^k F-x_0P_1^y\in P_0^B O[[v]],
                 \qquad F\in O[[v]] .
\]
The letter idempotent is implicit here. Put $Q=P_0$. Translation
$v=Q-p$ identifies the two integral power series rings by
$p$-adic completeness. There are an integer $d\le k$ and
$x'\in O$ such that
\[
      p^d f=x'+g,\qquad g\in (Q^B/P_1^y)O[[v]],
      \qquad
      d\le\left\lceil\max(0,y-B/p+2+\log_p y)\right\rceil .
\]
If $py<B$ one can take $d=0$.

In fact $\alpha=p^{-k}x_0$ satisfies
$F\equiv\alpha P_1^y\bmod Q^B$ over the fraction field.
Thus the first $B$ coefficients of $\alpha P_1^y$ are integral.
At $Q=0$, $P_1=p\varepsilon+Q^p V$ with both
$\varepsilon,V$ integral units, as in T83A2. In the binomial
expansion the $j$th summand starts at degree $pj$ with a unit
times $\alpha\binom yj p^{y-j}$. Induction in $j$ shows that
each such scalar with $pj<B$ is integral. Subtract the earlier
integral summands before inspecting the coefficient of $Q^{pj}$.
For $py<B$ use $j=y$ to get $\alpha\in O$. Otherwise put
$j=\lfloor(B-1)/p\rfloor<y$. The factorial valuation or the
carry count gives $v_p\binom yj\le\lfloor\log_p y\rfloor$.
Since $j\ge B/p-1$, the denominator valuation of $\alpha$ is at
most $y-B/p+1+\log_p y$. These calculations normalize
$v_p(p)=1$ and permit arbitrary ramification. If $\alpha=0$
take $d=0$. Otherwise round its nonnegative denominator
valuation up to the next integer. This $d$ is at most $k$ and
satisfies the stated bound. With $x'=p^d\alpha$, division of
$p^dF-x'P_1^y$ by $Q^B$ is integral because its first $B$
coefficients vanish. This gives the claimed numerator of $g$.
% END D83L1

% BEGIN D83L2
\paragraph{D83L2. The finite integral lift in the small initial-index branch.}
In the notation of U83M1 put $h=b+\eta+S$, $L_0=\log_p(2h)$,
$U=h/p^2+(W_0+S+\eta)/p+L_0+1$ and $I$ equal to the native
band indicator with constant $320$. For this paragraph $b$ is
the lower thick size. For every complementary target basis
coordinate and every scalar of D83L1 the original
$\mathrm{LiftsAt}$ assertion holds at an integer exponent
$\le15U+1+\eta I$ whenever
\[
             y\le (s_c+S+\eta)/p+2U+\eta .                 \tag{D83.1}
\]
In particular this includes the native branch
$y\le\mathrm{AI}(F,c-1)+\eta$. To recall the estimate for this
last index, use the signed widths $s_i+S+\eta$ in R82W1.
The forward widths at $c,c+1$ are at most $B+\eta$ and
$W_0+S+\eta$; thereafter they are at most $H_*=h+\eta$.
The forward weighted bound of R82W3 for the least annihilator
in the original variable adds $c_H+3$ to each width, where
$c_H=\lceil1+\log_p(H_*+1)\rceil$. Round its quotient by
$q=p-1$ up to get $\mathrm{AI}$. With $D=W_0+S+\eta$,
$m_0=B+\eta$ and $c_H\le2+L_0$, summation gives
\[
  \mathrm{AI}\le1+\frac{m_0}{q}+\frac{D}{q^2}
       +\frac{H_*}{q^2(q-1)}+\frac{c_H+3}{q-1}.
\]
Here the geometric weights start at $1,q^{-1},q^{-2}$.
Since $p\ge11$, $m_0\le h+\eta$ and $H_*\le2h$, the difference
between this expression and $m_0/p$ is at most $2U$. Indeed
$m_0/(pq)\le1.1(h+\eta)/p^2$, the tail is at most
$0.27h/p^2$, and $D/q^2\le0.11D/p$. Also $\eta\le D$ and
$1+(5+L_0)/9\le2(1+L_0)$.
This proves the implication to (D83.1).

Take $p^d f=x'+g$ from D83L1. An integral lift of the constant
coordinate at the initial cut, multiplied by $p^R$, follows from
U83M1 by lifting its Wach basis vector and specializing at zero.
Frobenius commutes with this map. Apply the signed frame
construction U83M6 to the corresponding constant vector in the
source, multiplied by $x'$. Its generator difference has the
required divisibility since the constant is fixed. Its image has
the prescribed $p^{R+\sigma}x'$ coordinate modulo $v\Lambda^2$
after clearing denominators: the horizontal series for the two
lattices commute with the map by uniqueness from their constant
terms. This proof permits arbitrary initial integral coordinates,
whether or not the horizontal constant is itself in the source
Wach lattice at the terminal letter.

For $g$ take its single complementary coordinate in the target
frame. At the up level its pole order is at most $y$ and its
constant direction is in $\operatorname{Fil}^{y}$. At the down
level it vanishes to order at least $B$ and every direction
then passes the filtration test. The two distinguished factors
have disjoint sets of roots in characteristic zero. Moreover
$\Lambda g$ is integral. Thus U83M6 and the integral primitive
filtration test give an actual vector in the target Wach lattice
representing $p^\sigma g$; it has zero contribution on every
other letter. Apply U83M1 at the terminal cut to lift its
$p^R$ multiple through the integral source lattice. Such a
source vector has generator difference in $vN$ by tame
invariance and the trivial action on $N/vN$. Both witnesses
have only the permitted poles and their target images have
$\Lambda$-multiples in the translated span. Adding the witnesses
gives $\mathrm{LiftsAt}$ at $e=R+\sigma+d$ for $f$, with its
original target congruence. This construction uses no change
of coefficient ring.

Finally $R+\sigma\le6U$. If $py<B$ then $d=0$. Otherwise
(D83.1) and $\eta/p+\log_p y+1\le U$ give $d\le\eta+5U$.
When $I=1$ this is sufficient. When $I=0$, this very nonthick
letter is adjacent to the thick one and $py\ge B\ge b$.
Failure of the band inequality at this letter gives
$py-b>2(p\eta+pS+b/p+y+pL_0)$.
On the other hand (D83.1) and $s_c\le b+\eta$ give
$py-b\le S+2\eta+2pU+p\eta$.
Consequently $(p-2)\eta<2pU$, so $\eta<3U$. Hence
$e\le14U$ off the band as required. In the band the bound
$e\le11U+\eta$ suffices.
% END D83L2
